% Section 6 only. Paste this file into main.tex or use:
% \input{section6_genai.tex}
%
% Required packages in the main document:
% \usepackage{booktabs,longtable,tabularx,listings}
% Use XeLaTeX because the exact prompts contain Vietnamese Unicode text.
%
% The main document must also define:
% \newcommand{\code}[1]{\texttt{\detokenize{#1}}}

\section{GenAI Usage \& Code Refinement Log (Mandatory Appendix)}

\newcommand{\sourceA}{\textbf{Source A: Antigravity history}}
\newcommand{\sourceB}{\textbf{Source B: Codex development conversation}}

\subsection{Purpose and Scope}

This appendix consolidates the two AI-assisted development histories used by
the group:

\begin{enumerate}
  \item \sourceA;
  \item \sourceB, 
\end{enumerate}

The selected interactions are representative examples, not a replacement for
the complete raw records. The group will submit
the complete Antigravity and Codex conversations export as supplementary evidence. This
prevents the selected excerpts from being interpreted outside their original
context.

AI assistance was substantial. It supported initial scaffolding, debugging,
protocol refinement, test construction, explanation, and documentation. The
evidence below focuses on how the group made choices, questioned results,
requested corrections, compared the implementation with course theory, and
required executable verification. The group does not claim that all code was
written without AI; it claims ownership through review, testing,
understanding, and responsibility for the final system.

\subsection{Development Timeline Across Both Histories}

\begin{longtable}{p{0.08\textwidth}p{0.18\textwidth}p{0.30\textwidth}p{0.34\textwidth}}
\toprule
Phase & Source & Main activity & Evidence of team decision or verification\\
\midrule
P01 & A & Requirements and architecture & Chose Python, TCP control/UDP data, and a staged SAW-to-GBN strategy\\
P02 & A & Initial implementation & Required English comments/debug output and a structured project layout\\
P03 & A & First E2E execution & Real STOR/RETR/HASH run exposed a TID mismatch\\
P04 & A & Debugging and UX & Fixed TID agreement and client auto-connect after observed failures\\
P05 & A & Broader acceptance tests & Requested binary, PORT, concurrency, and packet-loss validation\\
P06 & A & Branch/test review & Compared branches, retained the stronger implementation, and corrected a sandbox test\\
P07 & B & Requirement audit & Required manual review, grouped unit tests, and real E2E evidence\\
P08 & B & Protocol upgrade & Required GBN, sliding window, Slow Start, AIMD, and adaptive RTO\\
P09 & B & Conceptual review & Questioned session identity, authentication keys, ports, and statefulness\\
P10 & B & Theory comparison & Identified that the controller differed from TCP Reno learned in class\\
P11 & B & Demonstration evidence & Expanded real evidence from three panels to eight complete scenarios\\
P12 & B & Documentation & Aligned the seven-section report with the final code and documented limitations\\
\bottomrule
\end{longtable}

\subsection{Selected Interaction Records}

\subsubsection{Interaction 1 -- Requirements, constraints, and staged strategy}

\sourceA

\textbf{Exact user input}

\begin{lstlisting}
Đọc yêu cầu đề bài cho rõ và hoàn thành toàn bộ yêu cầu giúp mình.
Yêu cầu: Ngôn ngữ Python.
Nếu trong bài có cần comment hay những dòng code print để debug thì đều
phải sử dụng tiếng Anh, không dùng tiếng Việt nhé.

Mình chưa hiểu về Project socket này, nên cũng chưa hiểu việc nên chạy/test
lại project thế nào. Bạn hướng dẫn và hỗ trợ mình hết nhé.
\end{lstlisting}

\textbf{Later decision recorded by the user}

\begin{lstlisting}
Mình nghĩ nên chọn Go-Back-N (GBN). Chiến thuật an toàn: Hãy làm
Stop-and-Wait trước để chắc chắn nắm trong tay số điểm của mục Reliable UDP.
Khi mọi thứ đã chạy mượt, bạn hãy nâng cấp code lên GBN.

Quản lý Users: Hardcode ở Giai đoạn 1, Đọc file ở Giai đoạn hoàn thiện.
\end{lstlisting}

\textbf{Selected raw AI output}

\begin{lstlisting}
Project yêu cầu làm một Hybrid FTP với 2 kênh tách biệt:
- Control Channel (TCP): truyền lệnh và mã phản hồi
- Data Channel (UDP): truyền dữ liệu file, với reliability layer tự xây dựng

Selective Repeat hay Go-Back-N cho UDP reliability?
Users: hardcode admin/admin123 hay đọc từ file config?
\end{lstlisting}

\textbf{Team review and refinement}

We did not begin with an unspecified ``FTP application.'' We fixed the
language, low-level socket constraint, English source-output convention,
TCP/UDP channel separation, and target directory structure. We selected GBN
as the eventual reliability algorithm but intentionally proposed a staged
implementation so that basic reliability could be executed before adding
window complexity. We also chose a simple credential source for the
coursework scope. This interaction established explicit engineering choices
that could later be checked against the final code.

\textbf{Final artefacts:} \code{common/}, \code{server/}, \code{client/},
\code{server_files/}, and \code{client_downloads/}.

\subsubsection{Interaction 2 -- TID mismatch discovered by real execution}

\sourceA

\textbf{Executed validation requested in the conversation}

\begin{lstlisting}
Start a server in a background thread.
Connect the real FTPClient.
Login with admin/admin123.
Upload with STOR.
Download with RETR.
Compare content byte-for-byte.
Run HASH SHA-256.
\end{lstlisting}

\textbf{Selected raw AI output}

\begin{lstlisting}
Root cause found: Server calls _next_tid() -> returns 1,
client uses tid=0 -> TID mismatch -> receiver ignores all packets.
Fix: both sides use tid=0 (safe because each transfer uses fresh sockets).
\end{lstlisting}

\textbf{Team review and refinement}

The first network test did not pass, which prevented an unsupported claim of
completion. The packet trace showed that the sender and receiver filtered on
different transfer IDs. The fix aligned both endpoints for the existing
fresh-socket-per-transfer design, after which the same STOR/RETR/HASH test was
rerun. The second run verified login, upload, download, byte equality,
SHA-256, navigation, listing, NOOP, and QUIT. This record is important because
the implementation was corrected from observed runtime evidence, not only
from generated source.

\textbf{Verification:} the rerun printed \code{ALL TESTS PASSED}; generated
test files were then removed.

\subsubsection{Interaction 3 -- Client usability failure and auto-connect}

\sourceA

\textbf{Exact user input}

\begin{lstlisting}
Vì sao mình vẫn không chạy được?
\end{lstlisting}

\textbf{Selected raw AI output}

\begin{lstlisting}
Vấn đề rất rõ: server đang chạy tốt, nhưng client chưa connect --
bạn cần gõ open trước khi login. Tuy nhiên, đây là UX không tốt.
Tôi sẽ fix để client tự động connect khi khởi động.
\end{lstlisting}

\textbf{Team review and refinement}

The failure was reproduced at the CLI boundary rather than dismissed as user
error. Inspection found that \code{main()} did not provide the default host,
so the client started without a control connection. The startup path was
changed to connect to \code{127.0.0.1:2121}. A follow-up validation simulated
the exact startup flow. A separate ``address already in use'' result was also
correctly distinguished from a client logic error because an existing server
already occupied the port.

\textbf{Final result:} users can start the client and proceed directly to
login when the server is available.

\subsubsection{Interaction 4 -- Broader test criteria requested by the user}

\sourceA

\textbf{Exact user input (abridged only for repeated spacing)}

\begin{lstlisting}
Test truyền File Nhị phân lớn...
Test Chế độ Active Mode (PORT)...
Test Đa luồng (Multi-client Concurrency)...
Test Giả lập Mất gói / Timeout (Loss Recovery)...

Những yêu cầu này thì của mình đã đạt chưa?
Project đến đây là đã xong đủ các yêu cầu chưa?
\end{lstlisting}

\textbf{Team review and refinement}

This request expanded acceptance beyond a small text-file happy path. It
required binary fragmentation/reassembly, the alternate PORT setup, several
simultaneous clients, and deliberate UDP loss. These categories later became
part of the permanent integration suite and demo evidence. The question
``đã đạt chưa?'' also required evidence instead of a simple completion
statement.

\textbf{Final verification scope:} binary byte equality, PASV/PORT,
concurrency, compression, ABOR, deterministic loss/retransmission, metadata,
directory operations, and hash comparison.

\subsubsection{Interaction 5 -- A failing sandbox test was investigated}

\sourceA

\textbf{Exact user input}

\begin{lstlisting}
Chỗ Test lõi hệ thống (Core tests): thì mình bị hiện là fail.
Bạn xem thử xem.
\end{lstlisting}

\textbf{Selected raw AI output}

\begin{lstlisting}
Lỗi nằm ở file test, chứ không phải lỗi hệ thống Server.
File test cũ dùng một đường dẫn tuyệt đối. Cơ chế bảo mật mới tự động
đưa đường dẫn tuyệt đối vào trong thư mục gốc.

Thay vì dùng đường dẫn tuyệt đối, test dùng đường dẫn lùi
../tên_thư_mục_bên_ngoài. Server phát hiện cố vượt khỏi sandbox
và trả về None.
\end{lstlisting}

\textbf{Team review and refinement}

The test was not weakened merely to become green. The intended security
property was clarified: the resolver must reject a path that escapes the
server root. An absolute path was not a valid escape attempt under the
application's virtual-path rules, so the test was replaced by a relative
traversal case. The corrected test directly exercised the intended
\code{commonpath} boundary and then passed.

\textbf{Final artefact:} the path-sandbox regression test in
\code{tests/test_core.py}.

\subsubsection{Interaction 6 -- Upgrade from Stop-and-Wait to Go-Back-N}

\sourceB

\textbf{Exact user input}

\begin{lstlisting}
Go-Back-N
+ Sliding Window
+ Slow Start đơn giản
+ AIMD
- Adaptive timeout
cải tiến prject the hg này, sau khi cải tiếng xng, viết ại test và run test
ần ựt the tiêu chí cũ, viết ại reprt
\end{lstlisting}

\textbf{Selected raw AI output}

\begin{lstlisting}
Sender manages:
base = oldest unacknowledged packet
next_seq = next unsent packet

ACK n means all packets from 0 through n were received.
On timeout, Go-Back-N retransmits the outstanding flight:
for seq in range(base, next_seq):
    send_packet(packets[seq], retransmit=True)
\end{lstlisting}

\textbf{Team review and refinement}

We separated three ideas that are easily confused:

\begin{itemize}
  \item Go-Back-N provides reliability, ordering, cumulative ACK, and
        retransmission;
  \item \code{max_window=16} is a fixed protocol/sender ceiling; and
  \item \code{cwnd} is a dynamic congestion-response variable.
\end{itemize}

The implementation introduced \code{base}/\code{next_seq}, cumulative ACK,
receiver rejection of out-of-order packets, window retransmission,
\code{cwnd}/\code{ssthresh}, simplified Slow Start and AIMD, SRTT/RTTVAR,
bounded RTO, and Karn's rule. Observable histories were retained so unit tests
could assert internal behaviour, not only final file equality.

\textbf{Final artefacts:} \code{common/rdt.py},
\code{common/protocol.py}, updated tests, and updated protocol documentation.

\subsubsection{Interaction 7 -- Authentication and session identity}

\sourceB

\textbf{Exact prompts from the discussion}

\begin{lstlisting}
này cs hande việc giữ key để biết tài khản nàddg gin k

nhưng the kiến thưcx tui hc thì để biết packet này đến từ acc nà thì phải
dùng key đc tạ ử handshake mà

ủa vậy àm s server biết à sessin nà

ủa bthg cx z hả, 1 sessin chỉ đc nghe 1 scket?
\end{lstlisting}

\textbf{Selected raw AI output}

\begin{lstlisting}
Authentication is stateful because it belongs to the TCP Session.
Commands do not carry a login key; the accepted TCP socket identifies
the control session.

Each accepted TCP connection creates one independent Session object.
The session retains authenticated user, CWD, TYPE, MODE, PASV/PORT,
rename source, and cancellation state.
\end{lstlisting}

\textbf{Team review and refinement}

We compared the project's stateful FTP control connection with a previously
learned stateless token/key model. Source review confirmed that each
\code{accept()} creates \code{Session(conn, addr, sid)} and one thread. USER
selects a candidate account; PASS changes authentication state; subsequent
commands are associated through the same accepted TCP socket. We also
documented the limitation: socket/session identity is not cryptographic
identity, and the UDP transfer ID plus CRC is not an account-authentication
HMAC.

\textbf{Learning outcome:} we can distinguish stateful connection identity
from stateless bearer-token authentication and from transport encryption.

\subsubsection{Interaction 8 -- The group challenged the TCP comparison}

\sourceB

\textbf{Exact user input}

\begin{lstlisting}
cgestin này hình như hiw khác vs tui đc hc ử tcp, tại s ại àm khác
\end{lstlisting}

\textbf{Selected raw AI output}

\begin{lstlisting}
Project chỉ mô phỏng các ý tưởng chính của TCP trên UDP,
chứ không triển khai nguyên vẹn TCP Reno/Tahoe.

TCP thường có Slow Start, Congestion Avoidance,
Fast Retransmit và Fast Recovery.

TCP:     effective_window = min(cwnd, rwnd)
Project: effective_window = min(cwnd, max_window)
\end{lstlisting}

\textbf{Team review and refinement}

We noticed that the initial congestion-control description did not completely
match TCP Reno taught in class. We therefore changed the report's claim to
``simplified TCP-inspired application-layer congestion control.'' The project
has no receiver-advertised \code{rwnd}, three-duplicate-ACK Fast Retransmit,
Fast Recovery, SACK, or ECN. Timeout resets \code{cwnd} to one packet; NAK
halves it. CRC corruption is conservatively treated as loss even though
corruption does not necessarily imply router congestion.

This is direct evidence of critical use: an apparent theoretical mismatch was
questioned and the final claim was narrowed rather than exaggerated.

\subsubsection{Interaction 9 -- Four-layer completion audit}

\sourceB

\textbf{Exact user input}

\begin{lstlisting}
hãy liệt kê các yêu cầu và kiểm tra xem có những yêu cầu nào đã được hoàn
thiện, theo 4 yếu tố, check code manually, check unit test, phân nhóm check
test, check test end to end
\end{lstlisting}

\textbf{Selected raw AI output}

\begin{lstlisting}
Main functionality is complete and all 17 tests pass.
However, several details are only partial:
- MODE B has no separate block framing.
- HASH is available but not automatic after every transfer.
- UDP has no session HMAC.
- Authentication is plaintext and has no TLS.
- There is no per-file write lock.
- Congestion control is simpler than TCP.
\end{lstlisting}

\textbf{Team review and refinement}

We required four distinct evidence types: requirement matching, manual source
review, grouped component tests, and real end-to-end execution. All 28
approved command handlers were located. The test result was 16/16 core tests
plus 1/1 comprehensive integration scenario. Negative findings were retained
in the report so a passing suite would not be misrepresented as
production-grade security or complete RFC behaviour.

\subsubsection{Interaction 10 -- Real demo evidence, not hypothetical output}

\sourceB

\textbf{Exact user inputs}

\begin{lstlisting}
chạy ại dem và screen sht ch tui

chạy hết ại các dem này và chụp màn hình

k phải viết thành hứng dẫn mà à trình bày cách thực hiện dem và kết quả
\end{lstlisting}

\textbf{Selected raw AI output}

\begin{lstlisting}
DROP seq=0
-> timeout at base=0
-> retransmit outstanding flight
-> upload PASS
-> SHA-256 VERIFIED

Client 0: PASS
Client 1: PASS
Client 2: PASS
Concurrency: PASS
\end{lstlisting}

\textbf{Team review and refinement}

The first evidence generator covered three categories. We requested all
manual scenarios, so it was expanded to eight panels: control/authentication,
directory management, PASV integrity, ASCII/compression, PORT, extended file
operations, deterministic GBN loss recovery, and concurrency. Sequence zero
is dropped once so every run contains an actual timeout and retransmission.
Passwords are masked, the transcript is searched for \code{FAIL} and
\code{MISMATCH}, and generated files are removed.

\textbf{Final result:} eight inspected screenshots, no failure/mismatch line,
and clean client/server data directories.

\subsection{Combined Refinement and Debugging Register}

\begin{longtable}{p{0.06\textwidth}p{0.23\textwidth}p{0.28\textwidth}p{0.31\textwidth}}
\toprule
ID & Initial issue & Team review/refinement & Verification\\
\midrule
R01 & Initial design scope unclear & Fixed TCP/UDP split, Python/native sockets, staged reliability plan & Architecture and import checks\\
R02 & Server/client TID mismatch & Aligned TID for fresh per-transfer sockets & STOR/RETR/HASH rerun\\
R03 & Client started disconnected & Added default auto-connect path & Startup/login reproduction\\
R04 & Only small text happy path & Added binary, active, concurrent and loss scenarios & Integration tests\\
R05 & Sandbox unit test asserted wrong path semantics & Replaced absolute path with real relative traversal & Core test passed\\
R06 & Stop-and-Wait limited efficiency & Added cumulative-ACK GBN window & Window/retransmission tests\\
R07 & Fixed timeout ignored RTT variation & Added SRTT, RTTVAR, bounded RTO and Karn's rule & Adaptive-RTO test\\
R08 & Fixed window mislabeled as congestion control & Added loss-responsive cwnd/ssthresh, Slow Start and AIMD & cwnd-history tests\\
R09 & Payload corruption & Added CRC-32 reject/NAK path & Corruption test\\
R10 & Random loss evidence could be flaky & Deterministically dropped first DATA packet & Repeatable timeout/retransmission\\
R11 & TCP receive assumed message boundaries & Added CRLF buffering and multiline response parsing & Command tests\\
R12 & Active-mode method/field and reuse problems & Renamed field and recreated/reissued endpoint setup & Repeated PORT E2E\\
R13 & Path traversal risk & Normalised paths and checked common root & Sandbox regression\\
R14 & Initial explanation sounded like TCP Reno & Documented missing rwnd/Fast Recovery/SACK/ECN & Theory/source comparison\\
R15 & Evidence covered only three areas & Expanded to eight real screenshots & No FAIL/MISMATCH\\
R16 & Documentation described obsolete SAW state & Rebuilt diagrams and terminology for final GBN code & Manual report/source audit\\
\bottomrule
\end{longtable}


\subsection{Known Limitations Retained After AI Review}

\begin{itemize}
  \item Credentials are basic in-memory/plaintext values; there is no TLS,
        password hashing, role model, or expiring token.
  \item UDP packets use CRC-32 for accidental corruption, not an HMAC for
        authenticity.
  \item The current transfer ID is not a cryptographic session key.
  \item MODE B is accepted but has no separate RFC-style block-framing layer.
  \item HASH verification is available but is not forced after every transfer.
  \item Transfer payloads are assembled in memory instead of streamed.
  \item There is no per-filename write lock for simultaneous writers.
  \item The controller is simplified and is not TCP Reno/CUBIC.
  \item There is no receiver-advertised window, Fast Retransmit, Fast
        Recovery, SACK, or ECN.
\end{itemize}

Retaining these findings is part of the critical audit. A complete test pass
proves the tested coursework behaviours, not production-grade security or
all possible FTP/TCP semantics.



\subsection{Final Provenance Statement}

We declare that AI contributed substantially to the development process. We
do not claim that the product was created without AI. We selected the final
requirements, reported failures, challenged theoretical inconsistencies,
requested and reviewed changes, executed the software, retained negative
findings, and accept responsibility for the final submission.

This appendix contains representative interactions from both histories. The
complete unedited \code{AI.md} and Codex conversation export are submitted
alongside it so that prompts and outputs can be verified in full context.
